% Abstract + MCQ "Strategy" subsection.
% Human (outline) prose is left unmarked; prose added beyond the outline is wrapped in \aiadd{}.
% Requires the preamble to define \aiadd, e.g. \newcommand{\aiadd}[1]{\textcolor{blue}{#1}}.
% Bibliography: refs/abstract_strategy.bib

% NOTE: The abstract was extracted into main.tex (acmart requires the abstract
% environment to appear before \maketitle). Only the MCQ ``Strategy'' content
% remains in this file.

\section{Multiple-Choice Benchmarks}

\subsection{Strategy}

\aiadd{A computerized adaptive test (CAT) chooses each question from how the
model has answered so far. It selects the item that reveals the most about the
model's ability, then stops once the estimate is precise
enough~\cite{wainer2000cat, vanderlinden2000cat, weiss1982}. This selection rests
on item response theory (IRT), which models the chance of a correct answer from
an item's difficulty and discrimination together with the model's latent
ability~\cite{hambleton1991}. Because the test skips items that add little, it
reaches a target precision with far fewer questions than a fixed
benchmark~\cite{maiapolo2024tiny}.}

Our work builds upon ATLAS~\cite{li2025atlas}, which showed the potential for
adaptive testing to dramatically reduce the number of items required for
predicting benchmark accuracy. Items calibrated using item response theory have
also been shown to provide information beyond
accuracy~\cite{lalor2016, rodriguez2021, vania2021}. Models that answer harder
questions correctly but miss some easy ones are appropriately ranked higher
under a CAT, even though a simple accuracy comparison cannot capture this
difference~\cite{li2025atlas}.

This paper shows the feasibility of recreating the ATLAS methodology on new
benchmarks, and it formalizes the requirements for using a CAT as a true
replacement for an MCQ benchmark. It also sets out where multidimensional IRT
(MIRT) helps and where it falls short for MCQ, a discussion largely missing
from ATLAS~\cite{li2025atlas}.
